%% ---------------------------------------------------------------- cover
\begin{titlepage}
\thispagestyle{empty}
\centering
\vspace*{0.17\textheight}

{\color{FluxRule}\rule{\linewidth}{1.4pt}}\par\vspace{1.5em}
{\Huge\bfseries\color{FluxBlue} Flux\par}
\vspace{0.6em}
{\LARGE A Decentralized Cloud for\\[0.3em] Autonomous Compute\par}
\vspace{0.9em}
{\large\color{FluxSlate} A short introduction\par}
\vspace{1.5em}
{\color{FluxRule}\rule{\linewidth}{1.4pt}}\par

\vspace{2.2em}
{\large Tadeas Kmenta\par}
\vspace{0.45em}
{\color{FluxSlate}\today\par}

\vspace{3em}
\begin{minipage}{0.88\linewidth}
\small\color{FluxSlate}\setlength{\parskip}{0.75em}
\setlength{\emergencystretch}{2em}\hyphenpenalty=2000
What \Flux{} is, as measured on 2026-09-03; how it works; what is changing in v9 and when; what it
is for; and what does not work yet. Every number here is taken from a longer technical paper that
cites source code for each claim about deployed behaviour, and nothing here goes beyond it.

Where this document describes something planned rather than running, the sentence says so.
Where a claim has a limit, the limit sits beside it in a grey box. That is the method, and it is
the reason the document can be trusted.
\end{minipage}

\vfill
{\footnotesize\color{FluxSlate}
Companion to \emph{Flux v9: A Decentralized Cloud for Autonomous Compute}.\par}
\vspace{1.2em}
\end{titlepage}

%% ---------------------------------------------------------------- in one page
\clearpage
\thispagestyle{empty}
\section*{In one page}
\addcontentsline{toc}{section}{In one page}

\textbf{What it is.} A decentralized cloud of \num{6489} independently operated, collateralised
machines running \num{1195} deployed applications, coordinated by its own blockchain. An operator
locks \flux{} to join; the chain decides who may host and who is paid. Nobody signs a contract with
anyone.

\textbf{What already runs on it.} Not a proposal. By census: web services, blockchain
infrastructure, and \num{27} scientific-simulation deployments doing real protein-folding work on
about \SI{9}{\percent} of the fleet's committed CPU. A general-purpose cloud with science on it, into
which AI is one arriving workload among several.

\textbf{The bet.} The next tenant is software, not a person. Every major cloud is built for a human
with a card and an account. \Flux{} is being rebuilt for a caller that is a program --- one that
provisions, pays, and releases capacity with no human anywhere. It can, for one narrow and decisive
reason: on \Flux{}, paying is a signed transaction, not a billing relationship. \emph{A program can
hold a key. It cannot hold a corporate account.}

\textbf{Why this fleet.} The thing that centralises frontier AI --- sharding one model across many
accelerators that must be wired together --- simply does not apply to a model that fits on one
device. And the models most people actually need fit on one device. A dispersed fleet with a great
deal of memory is a structurally sensible home for them.

\textbf{What changes.} In July 2027, at block \num{3787502}, operators stop being paid purely by new
issuance: a fifth of every application payment flows into a pool that pays them instead. Ten
parallel-asset chains become two. The shielded pool is retired. A collateral-weighted vote can remove
an application. Each has a date, and each is planned rather than shipped until it lands.

\textbf{What does not work yet.} Hosting an application earns a node nothing extra, by design --- so
the whole hosting guarantee rests on an eligibility gate that today is bound to neither the operator
nor the machine. Four owner identities together hold enough weight to remove any application. There
will be no chain-level privacy. One Foundation-run API is the sole source of where applications
live. The bridge is administered by four people from one organisation. None of this is hidden in the
long paper, and none is hidden here.

\clearpage
\pdfbookmark[0]{Contents}{toc}
\tableofcontents
\clearpage
